El cronograma del proyecto se proyectó sobre un horizonte de
veinticuatro semanas (\cref{fig:scdee-gantt})con 2 hitos: la conclusión
de la fase 0, que marca la disponibilidad de la arquitectura
base, y la conclusión de la fase 6, que marca el primer momento en que el sistema completa el
flujo crítico desde la entrega física hasta la corrección. La fase 12,
dedicada a seguridad y pruebas, no se ejecuta en bloque al final del
proyecto sino transversalmente, lo que permite
endurecer el sistema de forma incremental.

\begin{figure}[Diagrama de Gantt del proyecto]%
              {fig:scdee-gantt}%
              {Diagrama de Gantt del proyecto.}
  \begin{gantt}{1}{24}
    \taskbar[fase0]{Fase 0: Configuración base}{1}{2}\\
    \taskbar[fase1]{Fase 1: Autenticación y organizaciones}{2}{6}\\
    \taskbar[fase2]{Fase 2: Cursos académicos}{6}{7}\\
    \taskbar[fase3]{Fase 3: Asignaturas y membresías}{7}{9}\\
    \taskbar[fase4]{Fase 4: Exámenes y modelos}{8}{11}\\
    \taskbar[fase5]{Fase 5: Instancias y corrección}{10}{15}\\
    \taskbar[fase6]{Fase 6: Ingesta y reconocimiento}{13}{15}\\
    \taskbar[fase7]{Fase 7: Anotaciones}{15}{16}\\
    \taskbar[fase8]{Fase 8: Revisiones}{16}{18}\\
    \taskbar[fase9]{Fase 9: Notificaciones}{18}{19}\\
    \taskbar[fase10]{Fase 10: Exportación y búsqueda}{19}{21}\\
    \taskbar[fase11]{Fase 11: Mantenimiento}{21}{22}\\
    \taskbar[fase12]{Fase 12: Seguridad y pruebas}{17}{23}\\
    \taskbar[fase13]{Fase 13: SSO y despliegue}{23}{24}\\
    \milestone[h1]{Arquitectura base}{2}\\
    \milestone[h2]{Pipeline operativo}{15}
  \end{gantt}
\end{figure}
